\documentclass[conference,10pt]{IEEEtran}
\usepackage{booktabs}
\usepackage{array}
\usepackage{amsmath}
\usepackage{url}
\usepackage[hidelinks]{hyperref}
\hypersetup{pdftitle={What Are We Measuring? Benchmark Integrity and a Zone-Based Reframing for Emotion Recognition in Autistic Children},pdfauthor={Nikhil Mahesh}}
% URTC requires a minimum of 10 point font throughout, including tables and references.
\let\urtcnormalsize\normalsize
\renewcommand{\small}{\urtcnormalsize}
\renewcommand{\footnotesize}{\urtcnormalsize}
\renewcommand{\scriptsize}{\urtcnormalsize}
\renewcommand{\tiny}{\urtcnormalsize}

\title{What Are We Measuring? Benchmark Integrity and a\\
Zone-Based Reframing for Emotion Recognition in Autistic Children}
\author{\IEEEauthorblockN{Nikhil Mahesh}\IEEEauthorblockA{Pioneer Research}}

\begin{document}
\maketitle

\begin{abstract}
Automated emotion recognition is widely proposed as an assistive technology for
autistic children, and recent papers report 97.95--99.99\% accuracy on public
autism facial-expression benchmarks. Independently collected clinical corpora
report 40--78\% on the same task. We resolve this contradiction by auditing the
most widely used public archive. Re-acquiring it and verifying it byte-for-byte
against our earlier run, we find that it ships two overlapping dataset trees
sharing 45\% of the smaller tree's images, that 21.4\% of its files are
byte-identical duplicates of another file, and that 119 duplicate groups carry
mutually contradictory emotion labels on pixel-identical images. Holding
features, classifier, and label space fixed and varying only the split protocol,
six-category balanced accuracy spans 0.325 to 0.709 across five defensible
readings of the same archive. Duplicate leakage itself proves second-order,
worth 4.4 points, and no protocol we construct approaches the published figures;
the dominant effect is which subset one evaluates on. The label contradictions
are not uniform: 279 of 280 zone-level conflicts fall in negative affect and one
falls in positive/neutral. Annotators agree on valence and cannot agree within
negative valence, so the six-category schema is itself manufacturing the noise it
is scored against. Collapsing to three action-oriented Emotion Zones cuts
protocol-induced variance by 38\%, which argues for coarser labels and
multimodal context over finer face-only categories. We report chance-corrected
metrics throughout so a collapse to fewer classes is never credited for the
chance floor moving. This is a measurement audit, not a deployable classifier,
and not a claim about any child's internal state.
\end{abstract}

\begin{IEEEkeywords}
benchmark integrity, data leakage, label noise, affective computing, autism
\end{IEEEkeywords}

\section{Introduction}
Roughly 1 in 31 eight-year-olds in the United States is diagnosed with autism
spectrum disorder (ASD)~\cite{shaw2025}. Because emotional expression in
autistic people is frequently misread by non-autistic observers---the
\emph{double empathy problem}~\cite{nas2018}---automated emotion recognition is
repeatedly proposed as a communication aid for autistic children and their
caregivers.

The literature on this task is in an unusual state. Papers evaluating on public
autism facial-expression archives report accuracies of
97.95\%~\cite{hanumantharayappa2026}, 99.8\%, and 99.99\%. Papers evaluating on
independently collected, consented clinical corpora report far less: 40\% for a
commercial analyser on the EMBOA autism corpus~\cite{kiejdo2025}, 57.95--64.92\%
for four commercial classifiers on parent-reported autistic
children~\cite{kalantarian2020}, and 66.43\% for a classifier trained on
typically developing children and tested on autistic children~\cite{grossard2020}.
The gap between these two literatures is roughly 30 accuracy points, and it
tracks how the evaluation data were obtained rather than what the models are.

That correlation has recently become impossible to ignore. Between November 2025
and 2026, publishers retracted or flagged roughly 90 papers built on scraped
autism face datasets, including the source paper for the archive we audit
here~\cite{talaat2023retracted} and a facial emotion recognition
paper~\cite{gungor2025retracted}, on the grounds that the images were collected
from the internet without documented diagnosis, ethical oversight, or guardian
consent. That established the \emph{provenance} failure. It did not establish
whether the reported numbers were also, independently, measurement artifacts.

This paper supplies that missing measurement. Our contributions are:

\begin{enumerate}
\item A quantified integrity audit of the most widely used public autism
emotion archive, computed from the vendor's own directory layout, reporting
duplicate structure, cross-copy overlap, and cross-label duplicate groups
(Sec.~\ref{sec:archive}).
\item A leakage ablation that holds features, classifier, and label space fixed
and varies only the split protocol across five plausible readings of the
archive, separating protocol effects from signal (Sec.~\ref{sec:ablation}).
\item The observation that the label contradictions are concentrated almost
entirely in negative affect, which we argue is evidence that the categorical
schema is generating them (Sec.~\ref{sec:zones}).
\item A resulting argument for coarser action-oriented zones and multimodal
context over finer face-only categories, with chance-corrected metrics so the
collapse is not credited for the chance floor (Sec.~\ref{sec:beyond}).
\end{enumerate}

We are auditing an instrument, not building one. We make no claim about any
child's internal state, and we treat every label as an annotator judgment about
an image rather than as ground truth about a person.

\section{Related Work}
\textbf{Cross-population transfer is established.} Grossard \emph{et
al.}~\cite{grossard2020} ran the controlled experiment on lab-collected video
from 157 typically developing and 36 autistic children with subject-independent
cross-validation, obtaining 82.05\% within-population and 66.43\% when training
on typically developing children and testing on autistic children. Gaya-Morey
\emph{et al.}~\cite{gayamorey2026} repeated the design with twelve CNNs on an
adjacent population, observing roughly 90\% within-population against under 55\%
across. We do not claim this finding. We ask a different question: whether the
public benchmarks now used to measure it can support the measurement at all.

\textbf{Benchmark contamination is a mature field.} Barz and
Denzler~\cite{barz2020} found 3.3\% and 10\% duplicate rates in CIFAR-10 and
CIFAR-100 test sets and released cleaned splits. Northcutt \emph{et
al.}~\cite{northcutt2021} showed that label errors in test sets reorder model
rankings. Ramos \emph{et al.}~\cite{ramos2025} give the taxonomy we adopt:
\emph{hard} leakage for identical images across a split, \emph{soft} for near
duplicates, \emph{intra-dataset} for overlap within one dataset's own splits.
Adimoolam \emph{et al.}~\cite{adimoolam2026} name \emph{augmented duplicates}
and report 89\% duplication in a geospatial benchmark. Apicella \emph{et
al.}~\cite{apicella2025} formalise \emph{synthesis leakage}, where train and
test items derive from a common source under a generating transform. None of
this literature has examined a clinical affect dataset.

\textbf{Contamination in this specific archive has been asserted, not
measured.} The authors of the FERAC derivative~\cite{ferac2024} reported ``class
overlap and image duplication'' in this archive, removed duplicates, and dropped
two classes as unrecoverable. They did not report a detection method or
threshold, a duplicate rate, a cross-split leakage rate, or the effect of
cleaning. Two independent groups have now flagged the same archive; what remains
open is the quantity.

\textbf{Why the labels might be the problem.} Barrett \emph{et
al.}~\cite{barrett2019} challenge the premise that emotional state is readable
from facial configuration at all. Keating \emph{et al.}~\cite{keating2026},
using over 265 million motion-capture points across 4{,}896 expressions, find
autistic and non-autistic expressions differ in which muscles carry each
emotion and in kinematic smoothness, concluding the two groups may be
``essentially speaking a different language.'' Zane \emph{et
al.}~\cite{zane2019} report that autistic expressions are produced at
intensities non-autistic raters find ambiguous. If the expressions are ambiguous
under a non-autistic reading, forcing them into six mutually exclusive
categories should produce exactly the contradictions we measure.

\textbf{Ethical framing.} Nagy~\cite{nagy2024} argues that autism has been
mobilised by emotion-AI developers as a charismatic use case and a testbed.
We take that critique as a constraint on what this paper may claim.

\section{The Archive}
\label{sec:archive}
We audit the public Kaggle archive
\texttt{fatmamtalaat/autistic-children-emo\allowbreak tions-dr-fatma-m-talaat},
the dataset associated with~\cite{talaat2023retracted}. We re-acquired it and
verified that all 1{,}672 SHA-256 digests recorded in our earlier manifests are
present in the fresh download, so the audit and the model results in
Sec.~\ref{sec:ablation} describe the same bytes. All statistics below are
computed from the vendor's own directory layout, before any re-splitting: this
is what a downstream user receives.

\textbf{The archive ships two overlapping datasets.} It contains two top-level
trees (Table~\ref{tab:trees}) with different class balances and different
train/test splits. They share 323 unique images, 45.0\% of the smaller tree.
This is a sufficient explanation for a puzzling feature of the literature:
published descriptions of ``the Talaat dataset'' report its size variously as
348, 830, 1{,}603, and ``2500+'' images. Papers are describing whichever tree
they loaded.

\begin{table}[t]
\caption{The two dataset trees shipped in one archive.}
\label{tab:trees}
\centering
\setlength{\tabcolsep}{4pt}
\begin{tabular}{lrrrr}
\toprule
Tree & Files & Unique & Train & Test\\
\midrule
Balanced copy & 1{,}425 & 1{,}380 & 1{,}199 & 226\\
Imbalanced copy & 833 & 717 & 758 & 75\\
\midrule
\multicolumn{5}{l}{Shared unique images: 323 (45.0\% of smaller tree)}\\
\bottomrule
\end{tabular}
\end{table}

\textbf{One file in five is a duplicate.} Of 2{,}258 labelled images, only
1{,}774 are distinct by SHA-256; 484 files (21.4\%) are byte-identical copies of
another file. Duplicate groups reach size 5. Perceptual hashing at Hamming
distance 4 adds 163 further near-duplicate components. A filename-suffix pattern
(\texttt{\_aug\_N}) identifies 1{,}020 augmentation families, 335 of them larger
than one and the largest containing 9 members, so a substantial part of the
archive is synthetic variation on a smaller set of source images. This is
synthesis leakage in the sense of~\cite{apicella2025}, and unusually it is
\emph{announced in the filenames}: it was trivially detectable and shipped
regardless.

\textbf{Pixel-identical images carry contradictory labels.} Under the vendor's
six categories, 119 byte-identical duplicate groups (327 files) contain members
filed under more than one emotion. This is a stronger and cheaper claim than
label-error detection in the sense of~\cite{northcutt2021}: establishing it
requires no external ground truth and no annotator panel, only a hash. It is an
internal inconsistency of the artifact.

\section{How Much of the Number Is the Protocol?}
\label{sec:ablation}
Contamination matters only if it moves results. We therefore hold the features,
the classifier, and the label space fixed and vary \emph{only} how train and
test are separated. Features are frozen \texttt{ViT-B/16}
embeddings~\cite{dosovitskiy2021}; the head is a one-hidden-layer MLP with
epoch-wise early stopping; every protocol is run over ten seeds~\cite{bouthillier2021}.
The five protocols are all defensible readings of the archive: the two trees'
own shipped splits; training on one tree and testing on the other; a pooled
random split that ignores duplicates, which is what an \texttt{ImageFolder} plus
\texttt{train\_test\_split} pipeline produces; and a group-aware split over
connected components of the exact-duplicate, near-duplicate, and
augmentation-family graph~\cite{joeres2025,saeb2017}.

\begin{table}[t]
\caption{Same features, same classifier, same labels; only the split protocol
changes. ``Group leak'' is the share of test images whose duplicate or
augmentation family also appears in training. Ten seeds.}
\label{tab:ablation}
\centering
\setlength{\tabcolsep}{3pt}
\begin{tabular}{lrrrr}
\toprule
Protocol & Bal.\ acc. & Acc. & Group leak & $n_{\text{test}}$\\
\midrule
\multicolumn{5}{l}{\emph{Six categories}}\\
Vendor split, imbalanced tree & 0.325 & 0.661 & 26.7\% & 75\\
Vendor split, balanced tree & 0.521 & 0.521 & 30.5\% & 226\\
Train one tree, test other & 0.451 & 0.727 & 25.3\% & 75\\
Group-aware & 0.665 & 0.705 & 0.0\% & 172\\
Pooled random & 0.709 & 0.721 & 81.6\% & 339\\
\midrule
\multicolumn{5}{l}{\emph{Three zones}}\\
Vendor split, imbalanced tree & 0.497 & 0.751 & 26.1\% & 69\\
Vendor split, balanced tree & 0.562 & 0.627 & 33.3\% & 189\\
Train one tree, test other & 0.665 & 0.841 & 24.6\% & 69\\
Group-aware & 0.727 & 0.765 & 0.0\% & 180\\
Pooled random & 0.735 & 0.779 & 80.6\% & 293\\
\bottomrule
\end{tabular}
\end{table}

Table~\ref{tab:ablation} supports three claims, and refutes a fourth that we
had expected to make.

\textbf{The reported number is largely a choice.} Across five defensible
protocols on identical data, six-category balanced accuracy spans 0.325 to
0.709: a 38.5-point range attributable entirely to how the archive is read. A
reader given only a single number cannot tell where in that band it sits.

\textbf{Metric choice compounds it.} Because the vendor's splits are severely
imbalanced, plain accuracy and balanced accuracy diverge sharply---0.661 versus
0.325 on the imbalanced tree's own test split, a 33.6-point gap on one protocol
with one model. Reporting accuracy on these splits is itself inflationary.

\textbf{Zones are markedly more protocol-stable.} The same five protocols span
23.8 points under the three-zone collapse against 38.5 under six categories, a
38\% reduction in protocol-induced variance. The coarser target is not merely
easier; it is less sensitive to a methodological choice the reader cannot see.

\textbf{But duplicate leakage alone does not explain inflated results.} The
pooled random split leaks the duplicate or augmentation family of 81.6\% of its
test images and still exceeds the group-aware split by only 4.4 points
(0.8 under zones). This is a negative result and we report it as one: in this
archive, leakage is real but second-order. Critically, \emph{no} protocol we
could construct---including the most contaminated---approaches the 97.95--99.99\%
figures in the literature. Under this standard recipe the published numbers are
not reachable by contamination alone, so they must originate elsewhere: in
augmentation applied before splitting in downstream derivatives, in test sets of
a few dozen images, or in selection among many reported runs. Our audit
constrains the explanation without completing it.

\section{Why the Labels Contradict}
\label{sec:zones}
The contradictions are not spread evenly across the label space. Mapping the six
categories onto three valence/arousal zones---zone 1 (natural, joy), zone 2
(anger, fear), zone 3 (sadness)---and recounting gives Table~\ref{tab:conflicts}.

\begin{table}[t]
\caption{Cross-label duplicate groups at two label granularities. Both views use
the same byte-identical groups.}
\label{tab:conflicts}
\centering
\setlength{\tabcolsep}{4pt}
\begin{tabular}{lrr}
\toprule
& Six categories & Three zones\\
\midrule
Conflicted duplicate groups & 119 & 103\\
Conflicted files & 327 & 280\\
\midrule
\multicolumn{3}{l}{\emph{Conflicted files by class}}\\
Positive / neutral & 5 & 1\\
Negative affect & 292 & 279\\
\bottomrule
\end{tabular}
\end{table}

Two things follow. First, collapsing to zones absorbs 16 of the 119 categorical
contradictions (13.4\%): those are cases where annotators disagreed between
anger and fear, or between labels that fall inside one zone, and the coarser
schema simply does not record the disagreement. Second, and far more strikingly,
the residue is almost entirely negative: 279 of 280 zone-level conflicted files
are in zones 2 and 3, against one in zone 1.

Annotators agree, essentially perfectly, on whether an autistic child's
expression is positive or negative. They cannot agree on which negative emotion
it is. That pattern is what~\cite{keating2026} and~\cite{zane2019} predict: if
autistic expressions are produced in a register non-autistic raters find
ambiguous, the ambiguity should concentrate where the categories are closest
together, which is inside negative affect.

This reframes the contradictions. They are not simply sloppiness in one archive.
A six-way forced choice over ambiguous expressions \emph{manufactures}
contradiction, because an annotator who perceives an ambiguous negative
expression must nonetheless pick one box. The label schema is generating the
noise the models are then scored against.

\section{Beyond Face-Only Categorical Recognition}
\label{sec:beyond}
Two design consequences follow, and one caution.

\textbf{Coarser, action-oriented labels.} If the reliable signal is valence and
the unreliable signal is the specific negative category, then a label space
built on the reliable boundary is better founded than one built across it. This
is the argument for Emotion Zones: classes defined by what a caregiver would
\emph{do}---de-escalate, engage, comfort---rather than by which Ekman category a
non-autistic annotator would select. Zones are not merely a convenience; on
these data they are the granularity at which annotation is reproducible.

\textbf{The caution: do not credit the chance floor.} A collapse from six
classes to three raises the chance baseline from 1/6 to 1/3, and comparing raw
balanced accuracies across the two spaces silently rewards that. On our earlier
data a commercial classifier moved from 28.17\% (six categories) to 45.12\%
(three zones), an apparent 16.95-point gain. Chance-corrected, the same
relabeling is worth 3.88 points. The zone argument must rest on construct
validity and annotation reliability, which the evidence above supports, not on a
discrimination gain that is mostly an artifact of counting. We therefore report
Cohen's $\kappa$ alongside balanced accuracy throughout.

\textbf{Context beyond the face.} Table~\ref{tab:models} gives our
leakage-controlled reference results on a group-aware split with 20 training
seeds and a cluster bootstrap over duplicate groups. Two numbers matter. A
neurotypically trained commercial classifier reaches 37.7\% balanced accuracy
with a 95\% interval of [33.1, 42.5] against a majority-class floor of 33.3\%,
and its recall on zone 3 is 5.8\%: on this test set it is close to
indistinguishable from predicting the majority class. A probe trained on the
same population reaches 77.8\%. The gap is not that the task is impossible; it
is that a face-only model trained on non-autistic faces is the wrong instrument.

\begin{table}[t]
\caption{Leakage-controlled reference results. Group-aware split, 251 test images
in 181 duplicate groups, 20 seeds, cluster bootstrap over groups.}
\label{tab:models}
\centering
\setlength{\tabcolsep}{3pt}
\begin{tabular}{lrlrr}
\toprule
Model & Bal.\ acc. & 95\% CI & $\kappa$ & Zone-3 rec.\\
\midrule
Majority class & 0.333 & --- & 0.000 & 0.000\\
DeepFace (NT-trained) & 0.377 & [0.331, 0.425] & 0.115 & 0.058\\
Linear probe & 0.731 & [0.669, 0.794] & 0.638 & 0.606\\
MLP probe & 0.778 & [0.725, 0.834] & 0.710 & 0.653\\
\bottomrule
\end{tabular}
\end{table}

This is where multimodal and vision-language approaches earn their place, and
the reason is a data argument rather than an accuracy argument. Training a
face-only autism-specific model requires autism-specific face data, which is
precisely the resource whose acquisition produced the retraction wave. A
vision-language model brings broad prior knowledge about scenes, activities, and
posture without requiring any autism-specific training corpus, and can supply
the environmental and interaction context that a cropped face cannot. Given
that our clean face-only ceiling is 77.8\% and that~\cite{barrett2019} argues
facial configuration underdetermines emotional state in principle, the
remaining headroom is unlikely to be found in better face classifiers. It is
more likely to be found in context---and in reporting a zone with an explanation
a caregiver can check, rather than a category with a confidence score.

\section{Limitations and Ethics}
This is an audit of one archive and does not establish prevalence across the
field, though~\cite{ferac2024} reports consistent qualitative findings. Our
duplicate detection is exact-hash and perceptual-hash based; it will miss
re-encoded or cropped duplicates, so our rates are lower bounds. Cluster
bootstrap intervals are percentile intervals, which can under-cover~\cite{anglin2026}.
We use no subject identities because the archive supplies none, so our groups
are duplicate and augmentation families, not people; a subject-disjoint split
would likely be stricter still.

We report only aggregate statistics. No image, filename, hash, or per-image
prediction is released. We make no claim that any depicted child is autistic,
that any label reflects a felt emotion, or that any system built on this
analysis should be deployed. The archive's provenance is disputed and its
license is unstated; we redistribute nothing and cite the source so the audit
can be reproduced from it. Under the EU AI Act, emotion inference in education
is prohibited except under a medical exemption whose justification explicitly
references autism, which raises rather than lowers the evidentiary bar on the
benchmarks used to support such systems.

\section{Conclusion}
The most widely used public autism emotion archive ships two overlapping copies
sharing 45\% of the smaller, is 21.4\% duplicated, and contains 119 groups of
pixel-identical images with contradictory labels. Varying only the split
protocol moves six-category balanced accuracy across a 38.5-point band, so a
single reported number on this archive carries little information about a model.
Duplicate leakage is not the main culprit---it is worth 4.4 points---and no
protocol we could construct reaches the 97.95--99.99\% figures in the
literature, which we therefore cannot attribute to contamination alone.

The label contradictions concentrate almost entirely in negative affect, which
indicates that the six-category schema is manufacturing them rather than merely
recording them, and the three-zone collapse cuts protocol-induced variance by
38\%. Coarser action-oriented zones and multimodal context are better founded
responses than finer face-only categories. Any future benchmark in this area
should report a duplicate audit, a group-aware split, and chance-corrected
metrics as a condition of being believed---and, given how this archive was
assembled, should not be built by scraping at all.

\section*{Acknowledgment}
% AI USE DISCLOSURE -- EXACT WORDING PENDING FROM AUTHOR

\bibliographystyle{IEEEtran}
\bibliography{refs}
\end{document}
